\documentclass[12pt,a4paper]{article}
% Preâmbulo institucional comum — DEE/IFMA
\usepackage[utf8]{inputenc}
\usepackage[T1]{fontenc}
\usepackage[brazil]{babel}
\usepackage{lmodern}
\usepackage{microtype}
\usepackage{geometry}
\usepackage{xcolor}
\usepackage{graphicx}
\usepackage{booktabs}
\usepackage{tabularx}
\usepackage{array}
\usepackage{enumitem}
\usepackage{hyperref}
\usepackage{lastpage}
\usepackage{fancyhdr}
\usepackage{setspace}
\usepackage{titlesec}
\usepackage{tcolorbox}
\usepackage{ragged2e}

\geometry{a4paper,top=2.2cm,bottom=2.2cm,left=2.5cm,right=2.5cm}
\definecolor{azulpetroleo}{HTML}{0F4C5C}
\definecolor{ambar}{HTML}{C8892D}
\definecolor{cinzaclaro}{HTML}{F2F5F7}
\definecolor{cinzaescuro}{HTML}{374151}
\hypersetup{colorlinks=true,linkcolor=azulpetroleo,urlcolor=azulpetroleo,citecolor=azulpetroleo}
\setlength{\parindent}{1.25cm}
\setlength{\parskip}{0.35em}
\onehalfspacing
\titleformat{\section}{\large\bfseries\color{azulpetroleo}}{}{0pt}{}
\titleformat{\subsection}{\normalsize\bfseries\color{azulpetroleo}}{}{0pt}{}
\pagestyle{fancy}
\fancyhf{}
\fancyhead[L]{\small Departamento de Eletroeletrônica — DEE}
\fancyhead[R]{\small IFMA — Campus São Luís/Monte Castelo}
\fancyfoot[C]{\small Página \thepage\ de \pageref{LastPage}}
\renewcommand{\headrulewidth}{0.4pt}
\renewcommand{\footrulewidth}{0.4pt}
\newcommand{\campo}[1]{\textcolor{ambar}{\textbf{[#1]}}}
\newcommand{\instituicao}{Instituto Federal de Educação, Ciência e Tecnologia do Maranhão}
\newcommand{\campus}{Campus São Luís/Monte Castelo}
\newcommand{\departamento}{Departamento de Eletroeletrônica — DEE}
\newcommand{\cidade}{São Luís — MA}
\newcommand{\cabecalhoInstitucional}{%
\begin{center}
{\Large\bfseries\color{azulpetroleo}\instituicao}\\[2mm]
{\large\bfseries\campus}\\[1mm]
{\normalsize\departamento}
\end{center}
\vspace{2mm}
\hrule
\vspace{4mm}
}
\newtcolorbox{destaque}{colback=cinzaclaro,colframe=azulpetroleo,boxrule=0.6pt,arc=2mm}

\begin{document}
\cabecalhoInstitucional

\noindent{\sffamily\bfseries OFÍCIO Nº \campo{número/ano} — DEE/MTC/IFMA}

\vspace{1.5mm}
\begin{flushright}
São Luís — MA, \campo{data}.
\end{flushright}

\vspace{4mm}
\noindent À\\
\textbf{\campo{NOME DA EMPRESA}}\\
A/C: \campo{setor ou responsável}\\
\campo{cidade/UF}

\vspace{5mm}
\noindent\textbf{Assunto:} Solicitação de apoio institucional para modernização dos laboratórios do Departamento de Eletroeletrônica do IFMA.

\vspace{4mm}
Prezados(as) Senhores(as),

O Instituto Federal de Educação, Ciência e Tecnologia do Maranhão — IFMA, por intermédio do Departamento de Eletroeletrônica — DEE do Campus São Luís/Monte Castelo, apresenta o \textbf{Programa de Modernização Tecnológica dos Laboratórios do DEE} e consulta essa empresa quanto à possibilidade de apoio institucional por meio de doação de equipamentos, disponibilização de kits didáticos, licenças, treinamento, cooperação técnica ou outra modalidade de parceria compatível com sua política institucional.

O DEE atua na formação técnica e superior nas áreas de Eletrotécnica, Eletrônica, Automação e Engenharia Elétrica, além de desenvolver atividades de pesquisa, extensão, iniciação científica e trabalhos de conclusão de curso.

Considerando a reconhecida atuação da \campo{NOME DA EMPRESA} na área de \campo{área de atuação}, identificamos especial aderência entre o portfólio da empresa e as necessidades acadêmicas do Departamento.

\vspace{2mm}
\noindent{\sffamily\bfseries\color{cinzaescuro} Relação preliminar de equipamentos e soluções de interesse}
\vspace{1.5mm}

\renewcommand{\arraystretch}{1.18}
\begin{tabularx}{\textwidth}{>{\RaggedRight\arraybackslash}X >{\centering\arraybackslash}p{1.4cm} >{\RaggedRight\arraybackslash}X}
\toprule
\textbf{Equipamento/solução} & \textbf{Qtd.} & \textbf{Aplicação acadêmica}\\
\midrule
\campo{Item 1} & \campo{Qtd.} & \campo{Laboratório/uso}\\
\campo{Item 2} & \campo{Qtd.} & \campo{Laboratório/uso}\\
\campo{Item 3} & \campo{Qtd.} & \campo{Laboratório/uso}\\
\bottomrule
\end{tabularx}

\vspace{4mm}
A relação apresentada é preliminar e poderá ser ajustada à política institucional, ao portfólio e à disponibilidade da empresa. Também há interesse em equipamentos de demonstração, unidades de treinamento, equipamentos descontinuados comercialmente porém plenamente funcionais, excedentes de estoque e outros bens adequados ao uso educacional, desde que sua transferência seja juridicamente possível e tecnicamente recomendável.

O apoio poderá beneficiar diretamente estudantes e servidores vinculados aos cursos e projetos do Departamento, promovendo maior integração entre a formação acadêmica e as tecnologias utilizadas pelo setor produtivo.

Eventual doação ou parceria será formalizada pelos setores competentes do IFMA, observando-se integralmente a legislação e as normas administrativas aplicáveis.

Colocamo-nos à disposição para reunião técnica, apresentação do projeto, visita aos laboratórios e adequação da relação de equipamentos à política e à disponibilidade dessa empresa.

\vspace{7mm}
Atenciosamente,

\vspace{13mm}
\begin{center}
\textbf{Francisco dos Santos Viana}\\[-0.2mm]
{\small Chefe do Departamento de Eletroeletrônica — DEE}\\[-0.2mm]
{\small IFMA — Campus São Luís/Monte Castelo}\\[1mm]
{\footnotesize \href{mailto:francisco.viana@ifma.edu.br}{francisco.viana@ifma.edu.br} \quad | \quad (98) 98816-6263}
\end{center}
\end{document}
